\section{Proof strategy and entropy preliminaries}\label{sec:overview}
The two terms in \eqref{eq:closed-functional} describe different
obstructions. The removed mass $\nu$ may concentrate at gaps too
small to resolve at the current budget. Its contribution is the
truncated inverse-gap cost $T_h$. In the remaining mass $\kappa$,
collective sign uncertainty and identifying the relevant scales
produce the cost-entropy term $h\sqrt{\rho\Gamma_\kappa}$.

\subsection{Organization of the proof}
The lower proof in Section~\ref{sec:matching-lower} first obtains
$R(B)\ge T_h/32$ by a posterior-martingale calculation and an
adaptive midpoint decomposition. It then filters the dyadic core
costs. Either the tail already controls the desired rate, or the
filter preserves a constant fraction of the core entropy. Grouping
and trimming the surviving bands produces conditional experiments
that satisfy the required rarity and saturation constraints.
Their lower bounds add through a common-pair argument.

The upper proof in Section~\ref{sec:direct-upper} uses the same
quantile split as a known design choice. A likelihood bank races
an ordinary capped sign classifier on observations of the same
coin. Coupled corrections at successive accuracies are replicated
with independent Poisson counts. The resulting finite multilevel
estimator has uniform expected work at most $B$. A homogeneous
kernel estimate controls mixed work across scales.
Section~\ref{sec:assembly-contract} assembles all branches and
tracks the constants in Theorem~\ref{thm:direct-main}.
Section~\ref{sec:comparison} compares $\Q$ with the selection
functional of a companion general-law manuscript.

Appendix~\ref{sec:uf-foundations} proves the classification, work
and allocation bounds. Appendix~\ref{sec:lower-foundations}
proves weighted sign testing, the activation gate and the annular
one-coin obstruction. Appendix~\ref{sec:transfer-foundations}
proves finite convex comparison, the common pair for all capped
local policies, the adaptive midpoint decomposition and additive
localization. No statistical selection theorem is assumed.

\subsection{Relation to prior methods}
Lee and Valiant~\cite{lv} study adaptive estimation for mixtures
of unknown coins and provide a midpoint decomposition for
fully adaptive observations. We reproduce the finite-prefix
decomposition needed here, including observable types and the
passage to uniform expected budgets. Their sample-complexity
theorem is not used as a substitute for the RMS statement above.

The weighted sign-testing argument follows the gap-entropy
mechanism of Chen, Li and Qiao~\cite{clq}, including their
Bernoulli extension. The specialized separated-gap lemma is
proved in Appendix~\ref{sec:lower-foundations}. The finite
convex comparison has methodological antecedents in
Juditsky and Nemirovski~\cite{jn} and Polyanskiy and Wu~\cite{pw};
its exact version and the simultaneous common-pair construction
are proved in Appendix~\ref{sec:transfer-foundations}.
We use Hoeffding's inequality~\cite{hoeffding} and Sion's convex
minimax theorem~\cite{sion}, together with the standard probability
and compactness facts indicated where applied. The resulting
quantile formula and all specialized reductions are established
within this manuscript. These attributions identify methodological
dependencies and do not assert a comprehensive novelty comparison.

Aronow, Kallus and Lopatto~\cite{akl} prove the gap-entropy
conjecture for fixed-confidence best-arm identification among
independent unit-variance Gaussian arms with means in $[0,1]$ and a
unique best arm $i^*$. For algorithms that identify the best arm with
probability at least $1-\delta$ on every such instance, in the range
of $\delta$ treated there, the optimal expected sample count averaged
over permutations of the arm labels is within constant factors of
$H\,(\log(1/\delta)+\Ent(\pi))$. Here $H=\sum_{i\ne i^*}\gamma_i^{-2}$,
$\gamma_i$ is the gap from arm $i$ to the best mean, and $\pi$ gives
the shares of $H$ in dyadic gap ranges. The core term
$\Gamma_\kappa$ has the same form. That problem identifies one of
finitely many labelled arms with a confidence requirement; here the
target is a population fraction under RMS loss and an expected
budget, with fresh unlabelled coins and unknown signs. We do not use
their result and do not claim that either result implies the other.

The companion manuscript~\cite{gl} states a characterization of
$R_{\mu,\rho}(B)$ for every law on $(0,1/2]$ by a functional
$\mathcal K$ that optimizes over submeasures of $\mu$; its proof has
not been externally refereed or formally verified.
Section~\ref{sec:comparison} shows that $\Q$ and $\mathcal K$ agree
within the factor $1+2\sqrt2$. Theorem~\ref{thm:direct-main} is
therefore the finite-law case of that stated characterization, here
with a closed functional and a self-contained proof that does not
invoke the theorem of~\cite{gl}. Both manuscripts were developed from
the same research notes: the two proofs share the observable upper
construction and most appendix lemmas and differ mainly in the
lower-bound reduction. The comparison itself uses neither statistical
theorem.

\subsection{Cost entropy and canonical band costs}\label{sec:entropy-preliminaries}
For a finite nonnegative vector $a$, put $A=\sum_j a_j$ and define
$\mathcal G(a)=A+\sum_{j:a_j>0}a_j\log(A/a_j)$, with
$\mathcal G(0)=0$. This continuous function is homogeneous,
coordinatewise increasing and concave. On the positive coordinates,
its derivative is $1+\log(A/a_j)$ and its Hessian is
$A^{-1}\mathbf1\mathbf1^\top-\operatorname{diag}(1/a_j)$.
Cauchy--Schwarz makes the latter negative semidefinite, and continuity
extends the inequalities to zero coordinates. Homogeneity and concavity
imply superadditivity, and binary entropy gives
\begin{equation}\label{eq:binary-mixture}
 \mathcal G(a+b)\le(1+\log2)[\mathcal G(a)+\mathcal G(b)]
                 \le2[\mathcal G(a)+\mathcal G(b)].
\end{equation}
Indeed, keep and then forget the label specifying which cost vector
generated a coordinate. The extra entropy is at most $\log2$ times
the total cost, which is at most $\mathcal G(a)+\mathcal G(b)$.

In the lower bound we use canonical costs
\[
 v_j=\kappa(I_j),\quad a_j=v_j/b_j^2,\quad A=\sum_j a_j,
 \quad F=\mathcal G(a),\quad \lambda=\sqrt{\rho/F}.
\]
The core has mass $1-\rho>0$, so $F>0$. Coordinatewise comparison yields
\begin{equation}\label{eq:canonical-core}
                      F\le\Gamma_\kappa\le4F.
\end{equation}
Only costs are rounded; the observation experiment retains its actual gaps.

